\chapter{Die zwei Kalibrierungen}

Die WDBT+ verwendet zwei unabhängige Kalibrierungen, die konsistent sein müssen.

\subsection{Kosmologische Kalibration}

Die kosmologische Kalibration basiert auf der gravitativen Rotverschiebung:

\begin{equation}
z(d) = \frac{4\pi G \rho_0 l_0^{3-D}}{D(D-1)c^2} \cdot d^{D-1}.
\end{equation}

Die Kalibration erfolgt durch SN Refsdal (\( z = 1,49 \)), deren Entfernung durch gravitative Linsenzeitverzögerung unabhängig bestimmt wurde. Daraus folgt:

\begin{equation}
\rho_0 \approx 5 \times 10^{-14} \, \text{kg/m}^3.
\end{equation}

\subsection{Atomare Kalibration}

Die atomare Kalibration basiert auf dem Lamb-Shift:

\begin{equation}
\Delta E_{\text{Lamb}}^{\text{WDBT+}} = \Delta E_{\text{Lamb}}^{\text{QED}} \cdot \left( \frac{a_0}{l_*} \right)^{D-3}.
\end{equation}

Die Kalibration mit dem experimentellen Lamb-Shift liefert:

\begin{equation}
l_* \approx a_0 \approx 5 \times 10^{-11} \, \text{m}.
\end{equation}

\subsection{Konsistenz}

Die Konsistenz wird durch die Vakuumenergiedichte hergestellt:

\begin{equation}
\rho_{\text{vac}} = \rho_0 \cdot \left( \frac{l_0}{l_*} \right)^{3-D}.
\end{equation}

Einsetzen der Werte:

\begin{equation}
\rho_{\text{vac}} \approx 5 \times 10^{-14} \cdot \left( \frac{1,8 \times 10^{-15}}{5 \times 10^{-11}} \right)^{0,296} \approx 10^{-9} \, \text{J/m}^3.
\end{equation}

Dies stimmt mit der beobachteten Vakuumenergiedichte überein.